\section{Related Work}
\label{sec:related_work}


\noindent \textbf{Vision language action models with memory.} 
Recent works have demonstrated that learned robot control policies, trained on large amounts of diverse robot experience, can lead to generalizable manipulation, e.g., in unseen environments~\citep{pi05, blackpi0, geminirobotics, lbmtri2025, nvidia2025gr00tn1openfoundation,ye2026worldactionmodelszeroshot}. Vision language action models (VLAs)~\citep{driess2023palm, rt22023arxiv, open_x_embodiment_rt_x_2023, octo_2023, Doshi24-crossformer, kim2024openvla, black2024pi_0, wen2024tinyvlafastdataefficientvisionlanguageaction, zhen20243dvla, belkhale2024minivla, szot2024multimodal, pertsch2025fast, geminirobotics2025, wen2025dexvla, bjorck2025gr00t, wu2026pragmatic, team2026gigabrain} are a popular approach for training such generalist policies, in which a pre-trained vision-language model is finetuned with robot experience. While many of today's state-of-the-art VLAs are trained without memory and act purely based on the current observation of the environment~\citep{kim2024openvla, blackpi0, pi05, geminirobotics, lbmtri2025, nvidia2025gr00tn1openfoundation}, a growing number of works have explored adding memory to policy training since it is a core requirement for solving a wide range of real-world tasks. While early works explore architectures with recurrent memory modules~\citep{rahmatizadeh2018vision, mandlekar2021matters}, more recent works that use transformer-based architectures simply pass a dense history of prior observations into the policy~\citep{shafiullah2022behavior, lee2024behavior, octo_2023}; computational and latency constraints make it challenging to scale such approaches to support very long-horizon memory. Some works have explored latent memory architectures~\citep{li2025cronusvla, shi2025memoryvla, fang2025sam2act, chung2025rethinking}, but only evaluated on short-horizon memory tasks. Others have proposed various heuristics to compress memory information, e.g., by relying on purely proprioceptive memory~\citep{zhang2025ta}, 2D point tracks~\citep{zheng2024tracevla, chen2025history}, by only retaining keyframes from prior timesteps~\citep{sridhar2025memer, wei2025cyclemanip, mark2026bpp}, or by representing memory in natural language~\citep{lin2025onetwovla}. %
A challenge for all these approaches is that it is hard to find a one-modality-fits-all solution for robot memory, and each individual representation will result in a compromise on capabilities: proprioception, point traces, or natural language alone, for example, lose precise spatial information about grasp angle and height that is necessary to correct a slipped grasp. Keyframes, on the other hand, need to aggressively sparsify the observation history to make inference computationally feasible in long-horizon tasks, but may lose the ability to estimate environment and robot dynamics, which often requires more densely sampled observations. 
In contrast to these prior works, we introduce a \emph{multi-modal} memory system that combines short-horizon, dense vision-based memory with long-horizon, language-based memory. This allows us to resolve partial observability, perform in-context adaptation, and solve long-horizon tasks, all without sacrificing efficiency. Finally, an orthogonal line of work proposes approaches for mitigating causal confusion~\citep{chi2025diffusion, zheng2024tracevla} in which a policy with memory erroneously learns to copy over prior actions, e.g., by introducing auxiliary objectives~\citep{torne2025learning}. While similar approaches \emph{could} be combined with our memory system, our experiments suggest that we can achieve high policy performance without such objectives, which may be attributed to our large-scale and diverse training data.




\noindent \textbf{Long-Context Models.} 
Outside of robotics, a large body of work has explored the training of long-context language and vision-language models \cite{liu2025comprehensive, tang2025video, ni2022expanding}. In particular, in the context of video processing, there is a rich literature on designing efficient encoders for long video inputs \cite{arnab2021vivit, bertasius2021space, li2024videomamba}. Our work leverages similar ideas to \cite{bertasius2021space, arnab2021vivit} for using sparse attention operations to efficiently process video inputs. Yet, in the context of robotics, latency constraints imposed by the real world require additional efficiency considerations: processing more than ten minutes of high-frequency, multi-frame video within a latency budget of a few hundred milliseconds is challenging, even on modern hardware. Thus, our work combines a short-horizon, video-based memory architecture with a significantly more compressed, language-based memory architecture for effective long-horizon context.
